\scanpage{175}
\mspage{4\deleted{\unsure{8}}9}

\deleted{\uline{Corollaire} \unsure{5.5.} Dans la définition d'une filtration
cardinale (SGA 4 I 9.12. p. 144), quand on
fait l'hypothèse de stabilité supplémentaire
d) sur les $\mathrm{Filt}^{\unsure{\pi|\varpi}}$, \inserted{$+$} dans la condition c) de \unsure{loc. cit.}
on peut supposer \deleted{que} \inserted{dans} la représentation de
$X$ comme $\underset{i\in I}{\underrightarrow{\mathrm{lim}}}\,X_i$ ($X_i\in\mathrm{Filt}^{\pi}$), $I$ grand
devant $\pi$, on ait \unsure{simplemt}, $I=\mathrm{Ss}_{\pi}(X)$}

\uline{Scholie} 5.5. Quand on fait \inserted{l'hyp. d'exactitude assez}
anodine que les $\underrightarrow{\mathrm{lim}}$ filtrantes de mono sont
mono, \uline{ou} qu'on se borne \unsure{aux cardinaux|au cardinal}
$\pi$ \unsure{tels|telle} que les objets de $\mathrm{Filt}^{\pi}$ soient
$\pi$-accessibles, on peut \uline{canoniser} la
représentation d'un objet $X$ comme $\underset{i\in I}{\underrightarrow{\mathrm{lim}}}\,X_i$,
avec $X_i\in\mathrm{Filt}^{\pi}$ et \deleted{\unsure{les}} $I$ grand devant $\pi$,
en prenant $I=\mathrm{Ss}_{\pi}(X)$, et le système
inductif correspondant. \deleted{(si \uline{lim} filtrante de mono est mono)}
On a de plus l'agréable propriété supplémentaire que
voici : toute partie filtrante \inserted{$J$} de $\unsure{I|II}$,
de cardinal $\leq\pi$, a une borne
supérieure $\deleted{\unsure{i}}\,i_J$, et (\deleted{si} \deleted{\unsure{l'on}} \unsure{suppose|supposons} que
$\unsure{\mathcal M|M}$ satisfait la condition \unsure{e)} page \unsure{46|45}, \unsure{bas page)|puisque l'|toujours)} on
aura
\[
 X_{i_J}=\underset{i\in J}{\underrightarrow{\mathrm{lim}}}\,X_i,
\]
en d'autres termes, on a une représentation
\deleted{\unsure{can}} de $X$ par un syst. inductif \inserted{(parfaitement)}
\uline{$\pi$-adapté}.
\inserted{(cf $\overline{\underline{\mathrm{XIII}}}$ p. 85, 86)}

% archive sheet number: 174
